\section{Introducción y Metodología de Análisis}

El estudio de la solicitación axil (tracción y compresión uniaxial) aborda el comportamiento resistente y deformacional de prismas mecánicos rectos sometidos a fuerzas cuya resultante actúa según la directriz longitudinal de la pieza.

El esquema secuencial de resolución de una estructura en Resistencia de Materiales sigue los siguientes pasos:
\begin{enumerate}
    \item \textbf{Modelización de la estructura:} Definición de la geometría mediante prismas mecánicos (barras, tirantes, vigas) y caracterización de las ligaduras y apoyos.
    \item \textbf{Cálculo de solicitaciones exteriores:} Identificación de las cargas aplicadas ($F$) y determinación de las reacciones de enlace mediante las ecuaciones de equilibrio estático de la estática clásica ($\sum \vec{F} = \vec{0}$, $\sum \vec{M} = \vec{0}$).
    \item \textbf{Cálculo de esfuerzos internos seccionales (2D):} Obtención de las leyes de esfuerzo axil ($N_x$), esfuerzo cortante ($Q_y$) y momento flector ($M_{fz}$):
          \begin{equation}
              \{ N_x(x), \quad Q_y(x), \quad M_{fz}(x) \}
          \end{equation}
          En este tema nos centraremos de forma específica en los problemas donde predomina o actúa exclusivamente el \textbf{esfuerzo axil} ($N_x \neq 0$, $Q_y = 0$, $M_{fz} = 0$).
    \item \textbf{Comprobación en las secciones críticas o más desfavorables:}
          \begin{itemize}
              \item \textbf{Criterio de Resistencia (Tensión Máxima):}
                    \begin{equation}
                        \sigma_{\text{máx},d} \le \sigma_{\text{adm,mat}} = f_d = \frac{f_k}{\gamma_m}
                    \end{equation}
              \item \textbf{Criterio de Rigidez / Deformación Excesiva (Deformación Límite):}
                    \begin{equation}
                        \varepsilon_{\text{máx}} \le \varepsilon_{\text{lím}} \quad \Longleftrightarrow \quad \Delta l \le \Delta l_{\text{adm}}
                    \end{equation}
          \end{itemize}
\end{enumerate}

\subsection{Valores Característicos ($k$) frente a Valores de Cálculo ($d$)}

En el cálculo estructural moderno según la normativa técnica (Eurocódigos, CTE, EAE):
\begin{itemize}
    \item Los valores reales esperados o de servicio se denominan \textbf{valores característicos} (subíndice $k$).
    \item Los valores ponderados de seguridad para dimensionamiento se denominan \textbf{valores de cálculo o de diseño} (subíndice $d$).
\end{itemize}

Para las acciones y cargas exteriores se introduce el \textbf{coeficiente de mayoración de cargas} $\gamma_F$ ($\gamma_F \ge 1.0$), de modo que:
\begin{equation}
    F_d = \gamma_F \cdot F_k
\end{equation}

Para las resistencias intrínsecas del material se introduce el \textbf{coeficiente de minoración del material} $\gamma_m$ ($\gamma_m \ge 1.0$), que depende de la naturaleza del material y del control de ejecución:
\begin{itemize}
    \item \textbf{Hormigón:} $\gamma_c \approx 1.50$ ($f_{cd} = f_{ck} / \gamma_c$).
    \item \textbf{Acero estructural:} $\gamma_s \approx 1.05\text{--}1.15$ ($f_{yd} = f_{yk} / \gamma_s$).
    \item \textbf{Madera estructural:} $\gamma_m \approx 1.30$ ($f_{md} = f_{mk} / \gamma_m$).
\end{itemize}

La condición de comprobación resistente se formula como:
\begin{equation}
    \sigma_{\text{máx},d} = \gamma_F \cdot \sigma_k \le \sigma_{\text{adm,mat}} = \frac{f_k}{\gamma_m}
\end{equation}

\begin{ImportanteBox}[title={Advertencia Crítica: Mayoración de Cargas vs. Cálculo de Deformaciones}]
    \textbf{¡OJO! Regla fundamental en dimensionamiento y examen:}
    \begin{itemize}
        \item \textbf{Para las TENSIONES ($\sigma$):} Se emplean los valores de los \textbf{esfuerzos mayorados} ($N_d = \gamma_F N_k$), comprobando frente a la resistencia de cálculo minorada ($\sigma_{\text{máx},d} \le f_d$).
        \item \textbf{Para las DEFORMACIONES y ALARGAMIENTOS ($\varepsilon, \Delta l$):} Se calculan \textbf{SIN MAYORAR}, es decir, en \textbf{Estado Límite de Servicio (ELS)} utilizando directamente los valores característicos de las cargas ($F_k, N_k$).
    \end{itemize}
    \textit{Estrategia recomendada:} Es preferible resolver la estática con las cargas características $F_k$ (obteniendo $N_k, \delta$), calcular los alargamientos con $N_k$, y mayorar al final multiplicando por $\gamma_F$ para verificar la tensión resistente $\sigma_d$.
\end{ImportanteBox}

\begin{figure}[H]
    \centering
    \adjustbox{max width=0.95\textwidth}{%
        \begin{tikzpicture}[>=Latex, node distance=1.8cm, auto]
            % Diagrama de flujo de cálculo con coeficientes de seguridad
            \node [draw, rectangle, rounded corners, fill=blue!10, text width=3.5cm, align=center, thick] (cargas) {\textbf{Cargas}\\\textbf{Características}\\$F_k$};
            \node [draw, rectangle, rounded corners, fill=yellow!15, text width=3.5cm, align=center, thick, right=1.8cm of cargas] (esfuerzos) {\textbf{Esfuerzos Axiles}\\$N_k, Q_k, M_k$};
            \node [draw, rectangle, rounded corners, fill=green!15, text width=3.8cm, align=center, thick, above right=0.8cm and 1.8cm of esfuerzos] (deform) {\textbf{Alargamientos (ELS)}\\$\Delta l = \frac{N_k l}{E \Omega} \le \Delta l_{\text{lím}}$\\\small(Sin mayorar)};
            \node [draw, rectangle, rounded corners, fill=red!10, text width=3.8cm, align=center, thick, below right=0.8cm and 1.8cm of esfuerzos] (tension) {\textbf{Tensiones de Cálculo (ELU)}\\$\sigma_d = \gamma_F \cdot \frac{N_k}{\Omega} \le \frac{f_k}{\gamma_m}$\\\small(Mayoradas con $\gamma_F$)};

            \draw [->, line width=1.3pt] (cargas) -- node[above, font=\footnotesize] {Equilibrio} (esfuerzos);
            \draw [->, line width=1.3pt] (esfuerzos) |- (deform);
            \draw [->, line width=1.3pt] (esfuerzos) |- node[pos=0.7, below, font=\footnotesize] {$\gamma_F$} (tension);
        \end{tikzpicture}%
    }
    \caption{Ruta de cálculo diferenciada para comprobaciones de rigidez (deformaciones sin mayorar) y resistencia (tensiones de cálculo mayoradas).}
    \label{fig:ruta_calculo_seguridad}
\end{figure}

%=============================================================================
\section{Tensiones en Tracción y Compresión Uniaxial}

\subsection{Esfuerzos Seccionales e Integrales de Tensión}

Consideremos un prisma mecánico en equilibrio estático cortado por una sección recta transversal $\Omega$ de normal exterior $\vec{n} = \vec{i}$ orientada según el eje longitudinal $x$. La reducción del sistema de fuerzas interiores sobre el centro de gravedad $G$ define el vector resultante $\vec{R}$ y el momento resultante $\vec{M}_R$:
\begin{align}
    \vec{R}   & = N_x \vec{i} + Q_y \vec{j} + Q_z \vec{k}     \\
    \vec{M}_R & = T \vec{i} + M_{fy} \vec{j} + M_{fz} \vec{k}
\end{align}

Las componentes de reducción se relacionan de forma exacta con la distribución continua de tensiones $\{\sigma_{nx}, \tau_{xy}, \tau_{xz}\}$ mediante las integrales de sección:
\begin{align}
    N_x    & = \iint_{\Omega} \sigma_{nx} \, d\Omega           \\
    Q_y    & = \iint_{\Omega} \tau_{xy} \, d\Omega             \\
    M_{fz} & = \iint_{\Omega} - y \cdot \sigma_{nx} \, d\Omega
\end{align}

En el caso específico de \textbf{solicitación axil pura}:
\begin{equation}
    N_x \neq 0, \qquad Q_y = 0, \qquad M_{fz} = 0, \qquad T = 0
\end{equation}

\begin{figure}[H]
    \centering
    \adjustbox{max width=0.85\textwidth}{%
        \begin{tikzpicture}[>=Latex, scale=1.1]
            % Cilindro / prisma mecánico tridimensional
            \draw[thick, fill=blue!5] (0,-1) arc (-90:90:0.4 and 1) -- (5,1) arc (90:-90:0.4 and 1) -- cycle;
            \draw[thick, fill=blue!15] (5,0) ellipse (0.4 and 1);
            \draw[dashed] (0,-1) arc (-90:90:0.4 and 1);
            \draw[thick] (0,-1) arc (270:90:0.4 and 1);

            % Eje longitudinal x
            \draw[dashdotted, gray!80, line width=1pt] (-0.8,0) -- (6.5,0) node[right, black] {$x$};

            % Sección intermedia de corte
            \draw[thick, fill=orange!25, fill opacity=0.8] (2.8,0) ellipse (0.4 and 1);
            \node[above=1.25cm, font=\small\bfseries, text=orange!80!black] at (2.8,0) {Sección transversal $\Omega$};

            % Distribución uniforme de tensiones sigma_nx
            \foreach \y in {-0.6, -0.3, 0.0, 0.3, 0.6} {
                    \draw[->, line width=1.1pt, red!85!black] (2.8,\y) -- ++(1.2,0);
                }
            \node[right, font=\small\bfseries, text=red!85!black] at (4.0, 0.3) {$\sigma_{nx}$};

            % Resultante N_x (con suficiente separación del eje)
            \draw[->, line width=2pt, blue!90!black] (5.4,0) -- ++(1.4,0) node[right, font=\bfseries] {$N_x$};
            \draw[<-, line width=2pt, blue!90!black] (-0.4,0) -- ++(-1.4,0) node[left, font=\bfseries] {$N_x$};
        \end{tikzpicture}%
    }
    \caption{Prisma mecánico bajo tracción pura y distribución continua de tensiones normales $\sigma_{nx}$ sobre la sección transversal.}
    \label{fig:prisma_axil_3d}
\end{figure}

\subsection{Cálculo de Tensiones y la Hipótesis de Bernoulli-Navier}

Para deducir la distribución analítica de tensiones normales $\sigma_{nx}$ a partir del esfuerzo axil $N_x$, se establecen las hipótesis fundamentales de la Resistencia de Materiales:
\begin{enumerate}
    \item \textbf{Material Elástico y Lineal (Ley de Hooke):}
          \begin{equation}
              \sigma_{nx} = E \cdot \varepsilon_x
          \end{equation}
          donde $E$ es el módulo de elasticidad longitudinal o módulo de Young.
    \item \textbf{Hipótesis de Bernoulli-Navier:}
          Las secciones transversales planas y perpendiculares a la directriz antes de la deformación \textbf{permanecen planas y perpendiculares} a la directriz deformada, sin distorsión apreciable en su propio plano.
\end{enumerate}

Como consecuencia de la hipótesis de Bernoulli, todas las fibras longitudinales experimentan exactamente el mismo alargamiento relativo:
\begin{equation}
    \varepsilon_x = \frac{\Delta l_x}{l_0} = \text{cte} \quad \forall (y,z) \in \Omega
\end{equation}

Sustituyendo en la ley de Hooke, la tensión normal resulta rigurosamente constante en todos los puntos de la sección transversal:
\begin{equation}
    \sigma_{nx} = E \cdot \varepsilon_x = \text{cte}
\end{equation}

Al ser $\sigma_{nx}$ constante en toda el área $\Omega$, puede extraerse fuera del signo integral del esfuerzo axil:
\begin{equation}
    N_x = \iint_{\Omega} \sigma_{nx} \, d\Omega = \sigma_{nx} \iint_{\Omega} d\Omega = \sigma_{nx} \cdot \Omega
\end{equation}

Despejando, se obtiene la expresión fundamental de la \textbf{tensión en tracción y compresión uniaxial}:
\begin{equation}
    \boxed{\sigma_{nx} = \frac{N_x}{\Omega}}
\end{equation}

\begin{figure}[H]
    \centering
    \adjustbox{max width=0.88\textwidth}{%
        \begin{tikzpicture}[>=Latex, scale=1.1]
            % Barra suspendida con peso y deformación
            \draw[line width=1.5pt] (-1, 3.5) -- (1, 3.5);
            \foreach \x in {-0.9, -0.6, -0.3, 0.0, 0.3, 0.6, 0.9} {
                    \draw[thin] (\x, 3.5) -- ++(-0.2, 0.3);
                }

            % Prisma original (línea discontinua)
            \draw[thick, fill=blue!5] (-0.8, 3.5) rectangle (0.8, 0.5);

            % Fibras longitudinales antes y después
            \draw[dashed, red!70!black, thick, fill=red!5] (-0.8, 0.5) rectangle (0.8, -0.3);
            \draw[<->, line width=1pt, red!85!black] (1.1, 0.5) -- node[right, font=\small\bfseries] {Alargamiento $\Delta l$} (1.1, -0.3);
            \draw[<->, line width=1pt] (-1.1, 3.5) -- node[left, font=\small] {$l_0$} (-1.1, 0.5);

            % Carga aplicada
            \draw[->, line width=2pt, blue!90!black] (0, -0.3) -- (0, -1.5) node[below, font=\bfseries] {$F = 10\text{ kN}$};

            % Diagrama de tensiones en la sección (lado derecho)
            \begin{scope}[shift={(4.5, 0.8)}]
                \draw[thick, fill=gray!15] (-0.8, -1.0) rectangle (0.8, 1.0);
                \draw[->, line width=1pt] (-1.5, 0) -- (2.5, 0) node[right] {$x$};
                \draw[->, line width=1pt] (0, -1.4) -- (0, 1.4) node[above] {$y$};

                % Bloque uniforme de tensiones
                \foreach \y in {-0.7, -0.35, 0.0, 0.35, 0.7} {
                        \draw[->, line width=1.2pt, red!85!black] (0.8, \y) -- ++(1.0, 0);
                    }
                \draw[line width=1.2pt, red!85!black] (1.8, -0.7) -- (1.8, 0.7);
                \node[right, font=\small\bfseries, text=red!85!black] at (1.8, -0.35) {$\sigma_{nx} = \frac{N}{\Omega} = \text{cte}$};
            \end{scope}
        \end{tikzpicture}%
    }
    \caption{Tracción simple: conservación de secciones planas (Bernoulli), alargamiento homogéneo de fibras y distribución tensional uniforme.}
    \label{fig:bernoulli_traccion}
\end{figure}

\subsection{Criterio de Signos para Esfuerzo Axil y Tensión Normal}

El signo del esfuerzo axil y de la tensión normal se corresponden unívocamente:
\begin{itemize}
    \item \textbf{Esfuerzo Axil Positivo ($N_x > 0$):} El vector resultante sale de la sección, produciendo \textbf{Tracción} ($\sigma_{nx} > 0$, alargamiento de fibras).
    \item \textbf{Esfuerzo Axil Negativo ($N_x < 0$):} El vector resultante entra hacia la sección, produciendo \textbf{Compresión} ($\sigma_{nx} < 0$, acortamiento de fibras).
\end{itemize}

%=============================================================================
\section{Deformaciones y Desplazamientos en Solicitación Axil}

\subsection{Respuesta Cinemática según el Tipo de Solicitación}

Cada esfuerzo interno induce un mecanismo específico de deformación en el medio continuo:
\begin{table}[H]
    \centering
    \caption{Correspondencia entre esfuerzos seccionales, tensiones y tipos de deformación.}
    \adjustbox{max width=\textwidth}{%
        \renewcommand{\arraystretch}{1.4}
        \begin{tabular}{c c l l}
            \toprule
            \textbf{Esfuerzo Interno}           & \textbf{Tensión Asociada} & \textbf{Tipo de Deformación}  & \textbf{Efecto Físico en las Fibras}                                 \\
            \midrule
            \textbf{Esfuerzo Axil ($N_x$)}      & $\sigma_{nx}$ (Normal)    & Cambio de tamaño longitudinal & Alargamiento / Acortamiento ($\varepsilon_x = \frac{\Delta l}{l_0}$) \\
            \textbf{Esfuerzo Cortante ($Q_y$)}  & $\tau_{xy}$ (Tangencial)  & Cambio de forma / distorsión  & Variación angular ($\gamma_{xy}$) sin cambio de volumen              \\
            \textbf{Momento Flector ($M_{fz}$)} & $\sigma_{nx}$ (Normal)    & Cambio de tamaño variable     & Tracción y compresión simultánea (curvatura seccional)               \\
            \bottomrule
        \end{tabular}%
    }
    \label{tab:esfuerzos_deformaciones}
\end{table}

\begin{figure}[H]
    \centering
    \adjustbox{max width=0.92\textwidth}{%
        \begin{tikzpicture}[>=Latex, scale=1.0]
            % 1. Axil: Alargamiento / Acortamiento
            \node[font=\bfseries] at (1.5, 1.8) {Esfuerzo Axil ($N_x \leftrightarrow \sigma_{nx}$)};
            \draw[thick, fill=blue!10] (0, 0.5) rectangle (3, 1.3);
            \draw[dashed, red, thick] (-0.3, 0.6) rectangle (3.3, 1.2);
            \draw[->, very thick, blue!80!black] (3, 0.9) -- (4, 0.9) node[right] {$N$};
            \draw[<-, very thick, blue!80!black] (-1, 0.9) -- (0, 0.9);
            \node[font=\footnotesize, align=center] at (1.5, 0.1) {Cambio de tamaño: $\varepsilon_x = \frac{\Delta l}{l_0}$};

            % 2. Cortante: Distorsión angular
            \begin{scope}[shift={(6.8, 0)}]
                \node[font=\bfseries] at (1.8, 1.8) {Esfuerzo Cortante ($Q_y \leftrightarrow \tau_{xy}$)};
                \draw[thick, fill=blue!10] (0, 0.5) rectangle (2.5, 1.5);
                \draw[dashed, red, thick] (0.4, 0.5) -- (2.9, 0.5) -- (2.5, 1.5) -- (0, 1.5) -- cycle;
                \draw[->, very thick, orange!90!black] (2.5, 1.5) -- (2.5, 2.3) node[above] {$Q$};
                \node[font=\footnotesize, align=center] at (1.5, 0.1) {Distorsión angular: $\gamma_{xy}$};
            \end{scope}

            % 3. Flector: Curvatura
            \begin{scope}[shift={(13.2, 0)}]
                \node[font=\bfseries] at (1.5, 1.8) {Momento Flector ($M_{fz} \leftrightarrow \sigma_{nx}$)};
                \draw[thick, fill=blue!10] (0, 0.5) rectangle (3, 1.3);
                \draw[dashed, red, thick] (0.2, 1.4) to[bend left=15] (2.8, 1.4) -- (3.0, 0.4) to[bend right=15] (0.0, 0.4) -- cycle;
                \draw[->, very thick, green!60!black] (3.4, 0.9) arc (0:120:0.5) node[right=4pt] {$M_f$};
                \node[font=\footnotesize, align=center] at (1.5, 0.1) {Curvatura: Tracción + Compresión};
            \end{scope}
        \end{tikzpicture}%
    }
    \caption{Comparativa cinemática de las deformaciones producidas por esfuerzo axil, cortante y momento flector.}
    \label{fig:mecanismos_deformacion}
\end{figure}

\subsection{Deducción Diferencial del Alargamiento Total}

Consideremos una barra elástica de longitud $l$ empotrada en su extremo superior $A$ y sometida a una carga axial $F$ en su extremo inferior $B$. Si aislamos un elemento diferencial de longitud inicial $dx$ comprendido entre dos secciones rectas transversales a cota $x$:
\begin{itemize}
    \item La sección a cota $x$ experimenta un desplazamiento longitudinal $u(x)$.
    \item La sección a cota $x + dx$ experimenta un desplazamiento $u(x) + du_x$.
    \item El alargamiento diferencial del elemento $dx$ es $du_x$.
\end{itemize}

Por definición, la deformación longitudinal unitaria $\varepsilon_x$ representa el cociente entre el alargamiento experimentado por el elemento y su longitud inicial:
\begin{equation}
    \varepsilon_x = \frac{du_x}{dx}
\end{equation}

Despejando el diferencial de desplazamiento y aplicando la ley de Hooke uniaxil ($\sigma_{nx} = E \varepsilon_x$):
\begin{equation}
    du_x = \varepsilon_x \, dx = \frac{\sigma_{nx}}{E} \, dx = \frac{N_x(x)}{E(x) \cdot \Omega(x)} \, dx
\end{equation}

El alargamiento o acortamiento total de un tramo continuo comprendido entre las secciones $A$ y $B$ resulta de la suma acumulada (integral) de las deformaciones de todos sus elementos diferenciales:
\begin{equation}
    \boxed{\Delta l_{AB} = \int_{x_A}^{x_B} du_x = \int_{x_A}^{x_B} \varepsilon_x \, dx = \int_{x_A}^{x_B} \frac{N_x(x)}{E(x) \cdot \Omega(x)} \, dx}
\end{equation}

\subsubsection{Caso Fundamental: Barra Homogénea de Sección y Esfuerzo Constantes}
Si a lo largo del tramo $AB$ el esfuerzo axil es constante ($N(x) = N_{AB} = \text{cte}$), la sección transversal no varía ($\Omega(x) = \Omega = \text{cte}$) y el material es uniforme ($E(x) = E = \text{cte}$), todos los términos salen fuera del signo integral:
\begin{equation}
    \Delta l_{AB} = \frac{N_{AB}}{E \cdot \Omega} \int_{x_A}^{x_B} dx = \frac{N_{AB} \cdot (x_B - x_A)}{E \cdot \Omega} = \boxed{\frac{N_{AB} \cdot l_{AB}}{E \cdot \Omega}}
\end{equation}

\begin{figure}[H]
    \centering
    \adjustbox{max width=0.88\textwidth}{%
        \begin{tikzpicture}[>=Latex, scale=1.1]
            % Techo empotrado
            \draw[line width=1.5pt] (-1.2, 4.0) -- (1.2, 4.0);
            \foreach \x in {-1.0, -0.7, -0.4, -0.1, 0.2, 0.5, 0.8, 1.1} {
                    \draw[thin] (\x, 4.0) -- ++(-0.2, 0.3);
                }
            \node[above right, font=\bfseries] at (0, 4.0) {$A$ (Origen $x_A=0$)};

            % Barra prismática vertical
            \draw[thick, fill=blue!8] (-0.6, 4.0) rectangle (0.6, 0.8);
            \node[left, font=\small] at (-0.7, 2.4) {$\Omega = 2\text{ cm}^2$};
            \draw[<->, line width=1pt] (1.0, 4.0) -- node[right, font=\small\bfseries] {$l$} (1.0, 0.8);

            % Elemento diferencial dx
            \draw[thick, fill=red!20] (-0.6, 2.8) rectangle (0.6, 2.4);
            \draw[<->, line width=0.9pt, red!85!black] (-0.9, 2.8) -- node[left, font=\footnotesize\bfseries] {$dx$} (-0.9, 2.4);
            \draw[<->, line width=0.9pt, blue!85!black] (-1.8, 2.4) -- node[left, font=\footnotesize\bfseries] {$du_x$} (-1.8, 2.15);

            % Barra alargada
            \draw[dashed, red!80!black, thick, fill=red!5] (-0.6, 0.8) rectangle (0.6, 0.1);
            \draw[<->, line width=1.1pt, red!85!black] (1.0, 0.8) -- node[right, font=\small\bfseries] {$\Delta l_{AB}$} (1.0, 0.1);
            \node[right, font=\bfseries] at (0.6, 0.1) {$B$ ($x_B=l$)};

            % Carga F en el extremo
            \draw[->, line width=2pt, blue!90!black] (0, 0.1) -- (0, -1.0) node[below, font=\bfseries] {$F = 10\text{ kN}$};

            % Explicación conceptual desplazada a la derecha
            \node[draw, rounded corners, fill=yellow!10, text width=6.2cm, right=3.2cm of {0, 2.2}, font=\small] {
                \textbf{Interpretación Física:}\\
                $\bullet$ Las secciones próximas a $A$ apenas se desplazan.\\
                $\bullet$ Las secciones inferiores acumulan la suma de todas las deformaciones de los elementos superiores.\\
                $\bullet$ El desplazamiento del extremo $B$ es el alargamiento total acumulado:
                \[ u_B = \Delta l_{AB} = \int_{0}^{l} \frac{N}{E\Omega} \, dx \]
            };
        \end{tikzpicture}%
    }
    \caption{Modelización diferencial del alargamiento de una barra prismática bajo carga axil.}
    \label{fig:alargamiento_diferencial}
\end{figure}

\subsection{Efecto Poisson y Deformaciones Transversales}

Cuando una pieza prismática se alarga longitudinalmente en la dirección del axil ($x$), se contrae transversalmente en las direcciones perpendiculares ($y, z$); recíprocamente, cuando se comprime longitudinalmente, se expande transversalmente.

La relación entre las deformaciones transversales unitarias y la deformación longitudinal viene gobernada por el \textbf{coeficiente de Poisson} $\nu$:
\begin{align}
    \varepsilon_x                             & = \frac{\Delta l_x}{l_{0x}} \\
    \varepsilon_y = \frac{\Delta l_y}{l_{0y}} & = - \nu \cdot \varepsilon_x \\
    \varepsilon_z = \frac{\Delta l_z}{l_{0z}} & = - \nu \cdot \varepsilon_x
\end{align}
El signo negativo garantiza la concordancia física: a tracción ($\varepsilon_x > 0$), las secciones transversales adelgazan ($\varepsilon_y, \varepsilon_z < 0$).

\begin{ImportanteBox}[title={¿Por qué se desprecia el efecto Poisson en el cálculo de vigas y barras?}]
    En elementos estructurales prismáticos tipo barra o viga, la longitud $l$ es varios órdenes de magnitud superior al canto o ancho de la sección ($l \gg b, h$). Dado que el coeficiente de Poisson es siempre inferior a la unidad ($\nu < 0.5$, típicamente $\nu \approx 0.20$ en hormigón y $\nu \approx 0.30$ en acero):
    \begin{itemize}
        \item Si una barra se alarga longitudinalmente $\Delta l_x = 1\text{ mm}$, su dimensión transversal apenas varía unas pocas milésimas de milímetro ($\Delta b \approx 0.001\text{ mm}$).
        \item Por consiguiente, en Resistencia de Materiales clásica se asume que \textbf{el área transversal $\Omega$ permanece constante} frente al cálculo de esfuerzos y tensiones sin incurrir en error apreciable.
    \end{itemize}
\end{ImportanteBox}

